\begin{theorem}[7.4.4]
\label{IV.2.s7.4.4-fr}
Bajo las hipótesis sobre $\boldsymbol P$ de \textup{(7.4.1)}, sea $A$ un
$\boldsymbol P$-anillo local noetheriano.
\begin{enumerate}
\item[\emph{(i)}] Para todo ideal primo $\mathfrak p$ de $A$,
$A_{\mathfrak p}$ es un $\boldsymbol P$-anillo.
\item[\emph{(ii)}] Supongamos además que $\boldsymbol P$ satisface
la condición $(\mathrm P_{\mathrm{IV}})$. Entonces, todo anillo
local que sea una $A$-álgebra esencialmente de tipo finito es un
$\boldsymbol P$-anillo.
\end{enumerate}
\end{theorem}

\begin{proof}
\begin{enumerate}
\item[(i)] Aplicando (7.4.3, (i)), basta comprobar que,
para todo ideal primo $\mathfrak p'$ de $A'=\widehat A$,
$A'_{\mathfrak p'}$ tiene fibras formales geométricamente
regulares. Ahora bien, esto se demostró en (\textbf 0, 22.3.3 y 22.5.8).

\oldpage[IV-2]{200}

\item[(ii)] Sea $B$ una $A$-álgebra esencialmente de tipo
finito que sea un anillo local. Si $\mathfrak p$ es el ideal primo
de $A$ sobre el cual está situado el ideal maximal de $B$, $B$
también es una $A_{\mathfrak p}$-álgebra esencialmente de tipo
finito (1.3.8); en virtud de (i), se puede suponer, pues, que
$\mathfrak p$ es el ideal maximal $\mathfrak m$ de $A$. Entonces
$B=C_{\mathfrak q}$, donde $C$ es una $A$-álgebra de tipo finito y
$\mathfrak q$ un ideal primo de $C$ situado sobre $\mathfrak m$.
Sea $\mathfrak n\supset\mathfrak q$ un ideal maximal de $C$
(necesariamente situado sobre $\mathfrak m$). Si se pone
$k=A/\mathfrak m$, $C/\mathfrak mC$ es una $k$-álgebra de tipo
finito; por tanto, el cuerpo residual $k'$ de $C$ en el ideal
maximal $\mathfrak n$ es \emph{finito} sobre $k$ (\textbf I, 6.4.2). En virtud
de (i), bastará probar que $C_{\mathfrak n}$ es un
$\boldsymbol P$-anillo, puesto que $C_{\mathfrak q}$ es un anillo
local de $C_{\mathfrak n}$ en un ideal primo. La cuestión se reduce así
a demostrar el
\end{enumerate}
\end{proof}
